\begin{tomtat}
	\begin{minipage}[t]{0.3\textwidth}
		\vspace{0pt}
		\begin{tomtat}
			\begin{center}
				\textbf{ĐỀ 29}
			\end{center}
			
		\end{tomtat}
	\end{minipage}
	\hfill
	\begin{minipage}[t]{0.5\textwidth}
		\vspace{0pt}
		\begin{tomtat}
			\begin{center}
				\includegraphics[scale=0.08]{images/29}
			\end{center}
		\end{tomtat}
	\end{minipage}
\end{tomtat}
\textbf{I. PHẦN TRẮC NGHIỆM}
\begin{ex}
	Điều kiện xác định của căn thức $\sqrt{2x-1}$ là:
	\choice{$x \ge \frac{1}{2}$}{$x \neq \frac{1}{2}$}{$x \le \frac{1}{2}$}{$x > \frac{1}{2}$}
	\loigiai{
		Căn thức $\sqrt{A}$ xác định khi $A \ge 0$.
		Do đó, $\sqrt{2x-1}$ xác định khi $2x-1 \ge 0$.
		$2x-1 \ge 0 \Leftrightarrow 2x \ge 1 \Leftrightarrow x \ge \frac{1}{2}$.
		\\ \textbf{Đáp Án A}
	}
\end{ex}

\begin{ex}
	Nghiệm của phương trình $3x-2=x+4$ là:
	\choice{$x=\frac{3}{2}$}{$x=1$}{$x=3$}{$x=-3$}
	\loigiai{
		Ta có: $3x-2 = x+4$
		$\Leftrightarrow 3x - x = 4+2$
		$\Leftrightarrow 2x = 6$
		$\Leftrightarrow x = \frac{6}{2}$
		$\Leftrightarrow x = 3$.
		\\ \textbf{Đáp Án C}
	}
\end{ex}

\begin{ex}
	Nghiệm của bất phương trình $5(x+1)-6 \le 0$ là:
	\choice{$x \le -\frac{1}{5}$}{$x \le \frac{1}{5}$}{$x \le \frac{6}{5}$}{$x \ge \frac{1}{5}$}
	\loigiai{
		Ta có: $5(x+1)-6 \le 0$
		$\Leftrightarrow 5x + 5 - 6 \le 0$
		$\Leftrightarrow 5x - 1 \le 0$
		$\Leftrightarrow 5x \le 1$
		$\Leftrightarrow x \le \frac{1}{5}$.
		\\ \textbf{Đáp Án B}
	}
\end{ex}

\begin{ex}
	Tính giá trị biểu thức $A=\sqrt{8}-\sqrt{(1-\sqrt{2})^2}$ ta được kết quả là:
	\choice{$\sqrt{2}-1$}{$3\sqrt{2}+1$}{$3\sqrt{2}-1$}{$\sqrt{2}+1$}
	\loigiai{
		Ta có: $A = \sqrt{8} - \sqrt{(1-\sqrt{2})^2}$
		$A = \sqrt{4 \cdot 2} - |1-\sqrt{2}|$.
		Vì $1 < \sqrt{2}$ (do $1^2=1$ và $(\sqrt{2})^2=2$), nên $1-\sqrt{2} < 0$.
		Do đó, $|1-\sqrt{2}| = -(1-\sqrt{2}) = \sqrt{2}-1$.
		$A = 2\sqrt{2} - (\sqrt{2}-1)$
		$A = 2\sqrt{2} - \sqrt{2} + 1$
		$A = \sqrt{2} + 1$.
		\\ \textbf{Đáp Án D}
	}
\end{ex}

\begin{ex}
	Hàm số $y=ax^2$ đi qua điểm $A(-1;2)$ thì giá trị của a là:
	\choice{$2$}{$-2$}{$\frac{1}{2}$}{$-\frac{1}{2}$}
	\loigiai{
		Hàm số $y=ax^2$ đi qua điểm $A(-1;2)$, nên tọa độ của A thỏa mãn phương trình hàm số:
		$2 = a(-1)^2$
		$2 = a \cdot 1$
		$a = 2$.
		\\ \textbf{Đáp Án A}
	}
\end{ex}

\begin{ex}
	Gieo 2 đồng xu cân đối và đồng chất 1 lần. Xác suất sao cho hai đồng xu xuất hiện mặt giống nhau là:
	\choice{$0$}{$\frac{1}{4}$}{$1$}{$\frac{1}{2}$}
	\loigiai{
		Gọi S là mặt sấp, N là mặt ngửa.
		Khi gieo 2 đồng xu, các kết quả có thể xảy ra (không gian mẫu $\Omega$) là:
		$\Omega = \{SS; SN; NS; NN\}$.
		Số phần tử của không gian mẫu là $n(\Omega) = 4$.
		Gọi A là biến cố “hai đồng xu xuất hiện mặt giống nhau”.
		Các kết quả thuận lợi cho biến cố A là: $A = \{SS; NN\}$.
		Số phần tử của biến cố A là $n(A) = 2$.
		Xác suất của biến cố A là: $P(A) = \frac{n(A)}{n(\Omega)} = \frac{2}{4} = \frac{1}{2}$.
		\\ \textbf{Đáp Án D}
	}
\end{ex}
\begin{ex}
	Bóng trên mặt đất của một cây dài 12 m. Biết các tia nắng mặt trời tạo với mặt đất một góc xấp xỉ bằng $40^\circ$. Chiều cao của cây (làm tròn đến chữ số thập phân thứ nhất) là:
	\choice{$14,3$ m}{$7,7$ m}{$9,2$ m}{$10,1$ m}
	\loigiai{
		Gọi chiều cao của cây là $h$ (mét).
		Gọi chiều dài bóng của cây trên mặt đất là $b = 12$ m.
		Góc tạo bởi tia nắng mặt trời và mặt đất là $\alpha = 40^\circ$.
		Ta có mối quan hệ lượng giác trong tam giác vuông (cây, bóng cây, tia nắng):
		$\tan(\alpha) = \frac{h}{b}$
		$h = b \cdot \tan(\alpha)$
		$h = 12 \cdot \tan(40^\circ)$.
		Sử dụng máy tính:
		$\tan(40^\circ) \approx 0,8390996$
		$h \approx 12 \cdot 0,8390996 \approx 10,069195...$
		Làm tròn kết quả đến chữ số thập phân thứ nhất, ta được $h \approx 10,1$ m.
		\\ \textbf{Đáp Án D}
	}
\end{ex}

\begin{ex}
	Tứ giác ABCD nội tiếp đường tròn và $\widehat{BAD} = 110^\circ$. Thì số đo $\widehat{BCD}$ bằng:
	\choice{$110^\circ$}{$70^\circ$}{$180^\circ$}{$360^\circ$}
	\loigiai{
		Trong một tứ giác nội tiếp, tổng hai góc đối diện bằng $180^\circ$.
		Do đó, $\widehat{BAD} + \widehat{BCD} = 180^\circ$.
		Thay $\widehat{BAD} = 110^\circ$, ta có:
		$110^\circ + \widehat{BCD} = 180^\circ$
		$\widehat{BCD} = 180^\circ - 110^\circ$
		$\widehat{BCD} = 70^\circ$.
		\\ \textbf{Đáp Án B}
	}
\end{ex}

\textbf{II. PHẦN TỰ LUẬN}
\begin{ex}
	Giải hệ phương trình sau:
	$$ \begin{cases} -x+2y=3 \\ 3x-y=1 \end{cases} $$
	\loigiai{
		Ta có hệ phương trình:
		$$ \begin{cases} -x+2y=3 \quad &(1) \\ 3x-y=1 \quad &(2) \end{cases} $$
		Từ phương trình (2), suy ra $y = 3x-1$.
		Thế $y = 3x-1$ vào phương trình (1):
		\begin{align*} -x + 2(3x-1) &= 3 \\ -x + 6x - 2 &= 3 \\ 5x &= 3+2 \\ 5x &= 5 \\ x &= 1 \end{align*}
		Thay $x=1$ vào $y = 3x-1$:
		\begin{align*} y &= 3(1) - 1 \\ &= 3 - 1 \\ &= 2 \end{align*}
		Vậy nghiệm của hệ phương trình là $(x;y) = (1;2)$.
	}
\end{ex}

\begin{ex}
	Rút gọn biểu thức: $P = \left( \frac{1}{\sqrt{a}+1} - \frac{1}{a+\sqrt{a}} \right) : \frac{\sqrt{a}-1}{2\sqrt{a}+2}$ với $a>0$ và $a \neq 1$.
	\loigiai{
		Điều kiện: $a>0, a \neq 1$.
		Ta có:
		\begin{align*} P &= \left( \frac{1}{\sqrt{a}+1} - \frac{1}{\sqrt{a}(\sqrt{a}+1)} \right) : \frac{\sqrt{a}-1}{2(\sqrt{a}+1)} \\ &= \left( \frac{\sqrt{a}}{\sqrt{a}(\sqrt{a}+1)} - \frac{1}{\sqrt{a}(\sqrt{a}+1)} \right) \cdot \frac{2(\sqrt{a}+1)}{\sqrt{a}-1} \\ &= \frac{\sqrt{a}-1}{\sqrt{a}(\sqrt{a}+1)} \cdot \frac{2(\sqrt{a}+1)}{\sqrt{a}-1} \end{align*}
		Vì $a>0$ và $a \neq 1$, nên $\sqrt{a}-1 \neq 0$ và $\sqrt{a}+1 \neq 0$.
		\begin{align*} P &= \frac{\cancel{(\sqrt{a}-1)}}{\sqrt{a}\cancel{(\sqrt{a}+1)}} \cdot \frac{2\cancel{(\sqrt{a}+1)}}{\cancel{(\sqrt{a}-1)}} \\ &= \frac{2}{\sqrt{a}} \end{align*}
		Vậy $P = \frac{2}{\sqrt{a}}$.
	}
\end{ex}

\begin{ex}
	Cho phương trình bậc hai $x^2-x-2=0$. Không giải phương trình, tính giá trị của biểu thức $A=x_1^2+x_2^2+\frac{1}{x_1}+\frac{1}{x_2}$.
	\loigiai{
		Phương trình $x^2-x-2=0$ có $a=1, b=-1, c=-2$.
		$\Delta = b^2-4ac = (-1)^2 - 4(1)(-2) = 1+8=9 > 0$.
		Phương trình có hai nghiệm phân biệt $x_1, x_2$.
		Theo hệ thức Vi-ét, ta có:
		$$ \begin{cases} x_1 + x_2 = -\frac{b}{a} = -\frac{-1}{1} = 1 \\ x_1 x_2 = \frac{c}{a} = \frac{-2}{1} = -2 \end{cases} $$
		Ta biến đổi biểu thức A:
		\begin{align*} A &= x_1^2+x_2^2+\frac{1}{x_1}+\frac{1}{x_2} \\ &= (x_1+x_2)^2 - 2x_1x_2 + \frac{x_1+x_2}{x_1x_2} \end{align*}
		Thay giá trị từ hệ thức Vi-ét vào biểu thức A:
		\begin{align*} A &= (1)^2 - 2(-2) + \frac{1}{-2} \\ &= 1 + 4 - \frac{1}{2} \\ &= 5 - \frac{1}{2} \\ &= \frac{10}{2} - \frac{1}{2} \\ &= \frac{9}{2} \end{align*}
		Vậy $A=\frac{9}{2}$.
	}
\end{ex}
\begin{ex}
	Một bó hoa gồm 3 bông hoa màu đỏ và 2 bông hoa màu vàng. Bạn A chọn ngẫu nhiên 2 bông hoa từ bó hoa đó. Xác định số phần tử của không gian mẫu và tính xác suất của biến cố “Trong 2 bông hoa được chọn, có đúng 1 bông hoa màu đỏ”.
	\loigiai{
		Ký hiệu hoa đỏ: $Đ_1; Đ_2; Đ_3$.
		Ký hiệu hoa vàng: $V_1; V_2$.
		Tổng số bông hoa là $3+2=5$.
		
		Chọn ngẫu nhiên 2 bông hoa từ bó hoa. Các cặp bông hoa có thể chọn (không gian mẫu $\Omega$):
		$\Omega = \{ (Đ_1;Đ_2); (Đ_1;Đ_3); (Đ_1;V_1); (Đ_1;V_2); (Đ_2;Đ_3); (Đ_2;V_1); (Đ_2;V_2); (Đ_3;V_1); (Đ_3;V_2); (V_1;V_2) \}$.
		Số phần tử của không gian mẫu là $n(\Omega) = 10$.
		
		Gọi B là biến cố: “Trong 2 bông hoa được chọn, có đúng 1 bông hoa màu đỏ”.
		Các kết quả thuận lợi cho biến cố B:
 $B = \{ (Đ_1;V_1); (Đ_1;V_2); (Đ_2;V_1); (Đ_2;V_2); (Đ_3;V_1); (Đ_3;V_2) \}$.
		Số phần tử của biến cố B là $n(B) = 3 \times 2 = 6$.
		
		Xác suất của biến cố B là:
		$P(B) = \frac{n(B)}{n(\Omega)} = \frac{6}{10} = \frac{3}{5}$.
	}
\end{ex}
\begin{ex}
	Nhân ngày sách và văn hóa đọc Việt Nam 21/4/2023, một nhà sách đã có chương trình giảm giá. Bạn An đến mua một quyển sách Toán và một quyển sách Tiếng Anh với tổng giá ghi trên hai quyển sách đó là 150.000 đồng. Nhưng quyển sách Toán được giảm giá 20\%, quyển sách Tiếng Anh được giảm giá 35\% nên An chỉ phải trả tổng số tiền là 108.000 đồng. Hỏi giá ghi trên mỗi quyển sách là bao nhiêu?
	\loigiai{
		Gọi giá ghi trên quyển sách Toán là $x$ (đồng).
		Gọi giá ghi trên quyển sách Tiếng Anh là $y$ (đồng).
		Điều kiện: $x > 0, y > 0$.
		
		Tổng giá ghi trên hai quyển sách là 150.000 đồng, nên ta có phương trình:
		$x + y = 150000 \quad (1)$
		
		Giá quyển sách Toán sau khi giảm giá 20\% là: $x - 0,20x = 0,8x$ (đồng).
		Giá quyển sách Tiếng Anh sau khi giảm giá 35\% là: $y - 0,35y = 0,65y$ (đồng).
		
		Tổng số tiền An phải trả là 108.000 đồng, nên ta có phương trình:
		$0,8x + 0,65y = 108000 \quad (2)$
		
		Từ (1) và (2), ta có hệ phương trình:
		$$ \begin{cases} x + y = 150000 \\ 0,8x + 0,65y = 108000 \end{cases} $$
		Giải hệ phương trình:
		Từ (1) $\Rightarrow y = 150000 - x$.
		Thế $y = 150000 - x$ vào (2):
		\begin{align*} 0,8x + 0,65(150000 - x) &= 108000 \\ 0,8x + 97500 - 0,65x &= 108000 \\ 0,15x &= 108000 - 97500 \\ 0,15x &= 10500 \\ x &= \frac{10500}{0,15} \\ x &= 70000 \end{align*}
		Thay $x = 70000$ vào $y = 150000 - x$:
		\begin{align*} y &= 150000 - 70000 \\ y &= 80000 \end{align*}
		(Thỏa mãn điều kiện $x > 0, y > 0$).
		
		Vậy, giá ghi trên quyển sách Toán là 70.000 đồng và giá ghi trên quyển sách Tiếng Anh là 80.000 đồng.
	}
\end{ex}
\begin{ex}
	Cho nửa đường tròn tâm O đường kính AB = 2R. Từ A vẽ tia tiếp tuyến Ax cùng phía với nửa đường tròn đối với AB. Từ điểm M trên Ax kẻ tiếp tuyến thứ hai MC với nửa đường tròn (C là tiếp điểm). MB cắt nửa đường tròn (O) tại D (D khác B).
	\begin{enumerate}
		\item[a)] Chứng minh A, M, C, O cùng nằm trên một đường tròn.
		\item[b)] Vẽ CH vuông góc với AB ($H \in AB$). Chứng minh rằng: $MA^2 = MD \cdot MB$ và MB đi qua trung điểm của CH.
	\end{enumerate}
\end{ex}

\begin{ex}
 Một công ty sản xuất dụng cụ thể thao nhận được một đơn đặt hàng sản xuất 8000 quả bóng tennis. Công ty này sở hữu một số máy móc, mỗi máy móc có thể sản xuất 30 quả bóng trong một giờ. Chi phí thiết lập các máy này là 200 nghìn đồng cho mỗi máy. Khi được thiết lập, hoạt động sản xuất sẽ hoàn toàn diễn ra tự động dưới sự giám sát. Số tiền phải trả cho người giám sát là 192 nghìn đồng một giờ (người này sẽ giám sát tất cả các máy hoạt động). Số máy móc công ty nên sử dụng là bao nhiêu để chi phí sản xuất là thấp nhất?

\end{ex}

\setcounterref{ex}{0}
